% Part: applied-modal-logics
% Chapter: epistemic-logic
% Section: introduction

\documentclass[../../../include/open-logic-section]{subfiles}

\begin{document}

\olfileid{aml}{el}{int}

\olsection{ভূমিকা}

মোডাল যুক্তিবিদ্যায় যেমন \emph{মোডাল বচন} ও তাদের মধ্যকার অনুসিদ্ধান্ত-সম্পর্ক আলোচনা করা হয়, জ্ঞানতাত্ত্বিক যুক্তিবিদ্যায় তেমনি \emph{জ্ঞানতাত্ত্বিক বচন} ও তাদের মধ্যকার অনুসিদ্ধান্ত-সম্পর্ক আলোচনা করা হয়। এখানে মোডাল অপারেটরগুলিকে সম্ভাবনা ও আবশ্যিকতার প্রতীক হিসেবে ব্যাখ্যা না করে, একস্থানিক সংযোজকগুলিকে জ্ঞানগত বা বিশ্বাসগত অর্থে ব্যাখ্যা করা হয়, যাতে জ্ঞান ও বিশ্বাসের মডেল তৈরি করা যায়। যেমন, নিচের দাবিগুলি প্রকাশ করতে চাইতে পারি:

\begin{enumerate}
\item রিচার্ড জানে যে ক্যালগারি আলবার্টায় অবস্থিত।
\item অড্রি মনে করে যে সোফায় একটি কুকুর থাকতে পারে।
\item রিচার্ড জানে যে অড্রি জানে, তার ক্লাস মঙ্গলবারে হয়।
\item সবাই জানে যে এক বছরে ১২ মাস।
\end{enumerate}
% BN-SRC-391: use the catalogued spelling of the 1962 book's author.
সমকালীন জ্ঞানতাত্ত্বিক যুক্তিবিদ্যার সূত্র প্রায়ই Jaakko Hintikka-র ১৯৬২ সালের \emph{Knowledge and Belief} বইয়ে খোঁজা হয়। বইটি এমন সময়ে লেখা, যখন যুক্তিবিদ্যায় সম্ভাব্য জগৎভিত্তিক অর্থতত্ত্বের ব্যবহার ক্রমে বাড়ছিল। বস্তুত, জ্ঞানতাত্ত্বিক যুক্তিবিদ্যায় অন্য মোডাল যুক্তিবিদ্যার বেশির ভাগ অর্থতাত্ত্বিক উপকরণই ব্যবহৃত হয়, তবে সেগুলির ব্যাখ্যা আলাদা। প্রধান পরিবর্তনটি হল \emph{অভিগম্যতা-সম্পর্ক} কী উপস্থাপন করে। জ্ঞানতাত্ত্বিক যুক্তিবিদ্যায় এই সম্পর্কগুলি কোনো ধরনের \emph{জ্ঞানগত সম্ভাবনা} উপস্থাপন করে। আমরা দেখব, জ্ঞানের যে ধারণাটির মডেল করছি, তার ওপরই অভিগম্যতা-সম্পর্কের কাঙ্ক্ষিত শর্তগুলি নির্ভর করে। আরও দেখব, জ্ঞানতাত্ত্বিক ব্যাখ্যা দিলে ফ্রেম-সঙ্গতি তত্ত্বে কী ঘটে। উপরের উদাহরণগুলিতে রিচার্ড ও অড্রি—দুই কর্তা এবং প্রত্যেকে কী জানে তার পারস্পরিক সম্পর্ক আছে। আমরা এমন বহু-কর্তা জ্ঞানতাত্ত্বিক যুক্তিবিদ্যা বিবেচনা করব, যেখানে এগুলি প্রকাশ করা যায়। বিপরীতে, এক-কর্তা জ্ঞানতাত্ত্বিক যুক্তিবিদ্যায় কেবল একজন কী জানে বা বিশ্বাস করে, তার কথাই বলা হয়।

\end{document}
